\clearpage
\phantomsection
\chapter{PHÂN TÍCH VÀ THIẾT KẾ HỆ THỐNG}

\section{Phân tích yêu cầu}

\subsection{Mục tiêu hệ thống}
Hệ thống được thiết kế để tự động hóa quá trình cân bằng cho một game Roguelike Deckbuilder. Các mục tiêu chính bao gồm:
\begin{itemize}
    \item Phát triển một game prototype hoàn chỉnh với các cơ chế gameplay cốt lõi.
    \item Xây dựng một AI agent có khả năng tự động chơi game để thu thập dữ liệu.
    \item Triển khai và so sánh hai phương pháp tối ưu hóa dựa trên thuật toán di truyền: một phương pháp tối ưu hóa trực tiếp với đơn mục tiêu và một phương pháp nhận biết cấu trúc với đa mục tiêu.
    \item Cung cấp công cụ trực quan hóa để phân tích quá trình tiến hóa và kết quả cân bằng.
\end{itemize}

\subsection{Kiến trúc tổng thể}
Hệ thống được thiết kế theo kiến trúc 3 tầng module hóa, đảm bảo sự tách biệt rõ ràng giữa các thành phần (Hình \ref{fig:architecture}).
\begin{itemize}
    \item \textbf{Core Game Layer:} Tầng thấp nhất, chứa toàn bộ logic của trò chơi, bao gồm hệ thống chiến đấu, quản lý bộ bài, bản đồ, và AI agent.
    \item \textbf{Simulation Layer:} Tầng trung gian, đóng vai trò cầu nối. Nó nhận các bộ tham số (genomes) từ tầng trên, áp dụng chúng vào game, chạy các lượt mô phỏng và tính toán các chỉ số fitness.
    \item \textbf{Optimization Layer:} Tầng cao nhất, chứa các thuật toán tối ưu hóa (GA và NSGA-II), quản lý quần thể và thực hiện các toán tử di truyền để tạo ra các thế hệ mới.
\end{itemize}
Ngoài ra, một \textbf{Data Layer} độc lập quản lý việc tải dữ liệu game từ các file JSON và ghi log kết quả.

\begin{figure}[H]
    \centering
    \includegraphics[width=0.9\textwidth]{SystemStructure.png}
    \caption{Kiến trúc tổng thể của hệ thống.}
    \label{fig:architecture}
\end{figure}

\section{Thiết kế chi tiết các module}

\subsection{Genome Representation}
Hai cách tiếp cận tối ưu hóa đòi hỏi hai cách biểu diễn genome khác nhau.

\subsubsection{BalanceGenome trong GA đơn mục tiêu}
Genome này biểu diễn trực tiếp từng tham số của game, tạo ra một không gian tìm kiếm lớn (~500 tham số). Nó bao gồm các dictionaries ánh xạ từ ID của đối tượng game đến hệ số nhân:
\begin{itemize}
    \item \texttt{CardActionScalars}: Hệ số nhân cho sát thương/giáp của từng hành động trên thẻ bài.
    \item \texttt{CardCostModifiers}: Mức thay đổi mana cost cho từng thẻ bài.
    \item \texttt{EnemyHealthScalars}: Hệ số nhân cho máu của từng loại kẻ địch.
    \item \texttt{EnemyActionValueScalars}: Hệ số nhân cho sát thương/giáp của từng hành động của kẻ địch.
\end{itemize}
Cách tiếp cận này cho phép điều chỉnh chi tiết nhưng dễ bị "curse of dimensionality".

\subsubsection{HierarchicalGenome trong NSGA-II đa mục tiêu}
Genome này gom nhóm các tham số theo một cấu trúc phân cấp, giảm không gian tìm kiếm xuống còn ~50 tham số. Các tham số được tổ chức theo các tầng:
\begin{itemize}
    \item \textbf{Global Multipliers:} Các hệ số nhân toàn cục cho sát thương, máu, giáp, giá vàng.
    \item \textbf{Progression Scaling:} Các hệ số điều chỉnh độ khó theo từng giai đoạn của game (Early/Mid/Late game).
    \item \textbf{Category Scaling:} Các hệ số nhân theo loại thẻ (Tấn công, Kỹ năng, Sức mạnh) và độ hiếm của thẻ/kẻ địch (1 đến 5 sao).
    \item \textbf{Room Distribution:} Các trọng số để điều chỉnh tần suất xuất hiện của các loại phòng (Quái, Tinh Anh, Sự kiện, Cửa hàng, Nghỉ ngơi).
\end{itemize}
Thiết kế này giúp thuật toán tìm kiếm hiệu quả hơn và đảm bảo các giải pháp có tính nhất quán về mặt thiết kế.

\begin{lstlisting}[language=C, caption={Ví dụ cách HierarchicalGenome tính toán scaled damage cho thẻ bài}, label={lst:hierarchical_scaling}]
public int GetScaledCardDamage(CardData card, int floor)
{
    // Bat dau voi global multiplier
    float scalar = GlobalDamageMultiplier;
    
    // Nhan voi category scalars (Attack/Skill/Power)
    scalar *= CardTypeScalars[card.Type];
    
    // Nhan voi rarity scalars (1-5 sao)
    scalar *= CardStarScalars[card.StarRating];
    
    // Nhan voi progression scaling (Early/Mid/Late game)
    scalar *= GetProgressionScalar(floor, ScalingType.Damage);
    
    // Ap dung override cho the cu the (neu co)
    if (CardDamageOverrides.TryGetValue(card.Id, out float ovr))
        scalar *= ovr;
    
    // Tinh toan gia tri cuoi cung
    var damageAction = card.Actions
        .FirstOrDefault(a => a.Type == ActionType.DealDamage);
    if (damageAction == null) return 0;
    
    return Math.Max(1, (int)Math.Round(damageAction.Value * scalar));
}
\end{lstlisting}

\subsection{AI Agent Design (HeuristicPlayerAI)}
AI agent được thiết kế với các heuristics dựa trên luật để mô phỏng một người chơi có trình độ trung bình. Logic ra quyết định của AI được chia thành các phần:
\begin{itemize}
    \item \textbf{Trong chiến đấu:}
        \begin{itemize}
            \item \textit{Ưu tiên sống sót:} Tính toán tổng sát thương sắp nhận và ưu tiên chơi các thẻ phòng thủ để chặn sát thương.
            \item \textit{Ưu tiên kết liễu:} Nếu có thể kết liễu kẻ địch trong lượt, AI sẽ ưu tiên các thẻ tấn công.
            \item \textit{Tối ưu hiệu quả:} Nếu không có nguy hiểm trước mắt, AI sẽ chọn chơi các thẻ tấn công hoặc kỹ năng dựa trên một hàm tính điểm, ưu tiên hiệu quả trên mỗi mana.
        \end{itemize}
    \item \textbf{Trên bản đồ:} AI chọn đường đi dựa trên điểm số của từng phòng, được tính toán dựa trên phần trăm máu hiện tại. Ví dụ, nó sẽ ưu tiên phòng Tinh Anh (Elite) khi máu cao để tìm phần thưởng tốt hơn, và ưu tiên phòng Nghỉ (Rest) khi máu thấp.
    \item \textbf{Lựa chọn phần thưởng:} AI ưu tiên các thẻ bài có độ hiếm cao hơn và cố gắng cân bằng giữa thẻ tấn công và phòng thủ trong bộ bài.
\end{itemize}

\subsection{Data Flow Diagram}
Hình \ref{fig:dataflow} minh họa luồng dữ liệu trong hệ thống, từ khi genome được tạo ra đến khi fitness được tính toán. Quá trình bao gồm: (1) Optimization Layer tạo genomes, (2) GenomeApplicator áp dụng genome vào GameData, (3) Simulation Runner chạy game với AI agent, (4) SimulationStats được thu thập, và (5) Fitness Evaluator tính toán fitness từ stats.

\begin{figure}[H]
    \centering
    \includegraphics[width=0.9\textwidth]{dataflow.png}
    \caption{Sơ đồ luồng dữ liệu trong hệ thống}
    \label{fig:dataflow}
\end{figure}

\subsection{Combat Sequence Diagram}
Hình \ref{fig:combat_sequence} mô tả trình tự tương tác giữa các thành phần trong một lượt chiến đấu điển hình. Mỗi lượt bắt đầu với việc rút bài, sau đó AI quyết định chơi thẻ nào, CombatManager xử lý actions, và kẻ địch phản ứng lại. Quá trình lặp lại cho đến khi một bên thua.

\begin{figure}[H]
    \centering
    \includegraphics[width=0.6\textwidth]{combat_sequence.png}
    \caption{Sequence diagram của một lượt chiến đấu}
    \label{fig:combat_sequence}
\end{figure}

\subsection{Class Diagram}
Hình \ref{fig:class_diagram} thể hiện các class chính và mối quan hệ giữa chúng. Các module được tổ chức thành các tầng rõ ràng, với Core Game Layer (GameController, CombatManager, DeckManager), AI Layer (HeuristicPlayerAI), và Optimization Layer (GeneticAlgorithm, NSGA2Optimizer).

\begin{figure}[H]
    \centering
    \includegraphics[width=0.6\textwidth]{class_diagram.png}
    \caption{Class diagram của các module chính}
    \label{fig:class_diagram}
\end{figure}

\subsection{Fitness Evaluation}
Đối với phương pháp structure-aware sử dụng NSGA-II, hệ thống tối ưu hóa đồng thời 3 mục tiêu:
\begin{enumerate}
    \item \textbf{Balance Score:} Đo lường độ cân bằng của game bằng cách so sánh các chỉ số thực tế (tỷ lệ thắng, HP khi thắng, tầng khi thua) với các giá trị mục tiêu (target). Balance Score được tính bằng:
    \begin{equation}
    f_{balance} = \sum_{i} w_i \cdot \left| \frac{\text{actual}_i - \text{target}_i}{\text{target}_i} \right|
    \end{equation}
    trong đó $w_i$ là trọng số quan trọng của metric thứ $i$.
    
    \item \textbf{Engagement Score:} Đánh giá sự hấp dẫn và đa dạng chiến thuật, dựa trên số lượng thẻ bài "khả dụng" (Viable Cards) và sự đa dạng trong các lối xây dựng bộ bài (Build Variety):
    \begin{equation}
    f_{engagement} = \alpha \cdot \frac{|\text{ViableCards}|}{|\text{TotalCards}|} + \beta \cdot \text{BuildDiversity}
    \end{equation}
    với $\alpha, \beta$ là các trọng số cân bằng. Càng nhiều thẻ bài khả dụng và càng đa dạng lối chơi thì Engagement Score càng cao.
    
    \item \textbf{Coherence Score:} Đo lường tính nhất quán trong thiết kế. Điểm số được tính bằng cách lấy 1 trừ đi tổng hình phạt (penalty):
    \begin{equation}
    f_{coherence} = 1 - \left( \text{TrapCardPenalty} + \lambda \cdot \sum_{j} (\text{scalar}_j - 1)^2 \right)
    \end{equation}
    trong đó TrapCardPenalty phạt nặng các "thẻ bài rác" (Trap Cards - thẻ được chọn nhưng có tỷ lệ thắng thấp), và $\lambda$ kiểm soát mức độ phạt cho các tham số xa giá trị chuẩn. Khi không có trap cards và các scalars gần 1.0, Coherence Score đạt giá trị tối đa.
\end{enumerate}
Đối với phương pháp GA truyền thống, ba chỉ số trên được kết hợp thành một hàm fitness duy nhất:
\begin{equation}
F_{total} = w_1 \cdot f_{balance} + w_2 \cdot f_{engagement} + w_3 \cdot f_{coherence}
\end{equation}